% Shared preamble for the second edition of Graviton Pressure Theory.
% A part's own file sets \gptpart, \gpttitle, \gptsubtitle and \gptdate, then % Shared preamble for the second edition of Graviton Pressure Theory.
% A part's own file sets \gptpart, \gpttitle, \gptsubtitle and \gptdate, then % Shared preamble for the second edition of Graviton Pressure Theory.
% A part's own file sets \gptpart, \gpttitle, \gptsubtitle and \gptdate, then % Shared preamble for the second edition of Graviton Pressure Theory.
% A part's own file sets \gptpart, \gpttitle, \gptsubtitle and \gptdate, then \input{../common/preamble}.
\documentclass[11pt,a4paper]{article}
\usepackage[T1]{fontenc}
\usepackage[utf8]{inputenc}
\usepackage{mathpazo}
\usepackage[scaled=0.92]{helvet}
\usepackage{amsmath,amssymb}
\usepackage[margin=2.5cm]{geometry}
\usepackage{microtype}
\usepackage{graphicx}
\usepackage{booktabs,array,tabularx}
\usepackage[font=small,labelfont=bf]{caption}
\usepackage{enumitem}
\usepackage[authoryear,round]{natbib}
\usepackage{titlesec}
\usepackage[titles]{tocloft}
\renewcommand{\cftsecfont}{\bfseries}\renewcommand{\cftsecpagefont}{\bfseries}
\renewcommand{\cftsecleader}{\hfill}\setlength{\cftsecnumwidth}{3.2em}
\usepackage{gpt-status}
\usepackage[hidelinks]{hyperref}
\setlength{\parskip}{4pt}
\setlength{\emergencystretch}{3em}
\setlength{\parindent}{0pt}
\setlist{itemsep=2pt,topsep=3pt}
\titleformat{\section}{\Large\bfseries}{\thesection}{1em}{}
\titleformat{\subsection}{\large\bfseries}{\thesubsection}{1em}{}
\renewcommand{\thesection}{\gptpart.\arabic{section}}
\renewcommand{\thetable}{\gptpart.\arabic{table}}
\renewcommand{\thefigure}{\gptpart.\arabic{figure}}
\renewcommand{\theequation}{\gptpart.\arabic{equation}}
\setcounter{tocdepth}{1}
\makeatletter\renewcommand{\@pnumwidth}{2.2em}\renewcommand{\@tocrmarg}{3em}
\makeatother
\bibliographystyle{plainnat}
\newcommand{\gptitle}{%
  \begin{titlepage}\centering
  \vspace*{3cm}
  {\Huge\bfseries Graviton Pressure Theory\par}\vspace{0.8cm}
  {\Large The Unified Framework\par}\vspace{0.2cm}
  {\large Individual Submission\par}\vspace{0.6cm}
  {\normalsize This document is part of a multi-part scientific framework\par}\vspace{2cm}
  {\normalsize Part \gptpart\ of 30\par}\vspace{0.8cm}
  {\huge\bfseries \gpttitle\par}\vspace{0.6cm}
  {\Large \gptsubtitle\par}\vspace{1.5cm}
  {\small\begin{minipage}{0.75\textwidth}\centering This submission is part of the broader Graviton Pressure Theory (GPT) project, a redefinition of gravitational interaction rooted in causal field dynamics and physical force transmission. Each document is designed to stand on its own; its full context emerges as part of the larger framework. For the complete series, see https://peopleoftheworld.net/graviton-pressure-theory/.\end{minipage}\par}
  \vfill
  {\large Author: Shareef Ali Rashada\par}\vspace{0.2cm}
  {\normalsize shareef@peopleoftheworld.net\par}\vspace{0.5cm}
  {\normalsize Second edition, \gptdate\par}
  \end{titlepage}
  \setcounter{page}{1}\tableofcontents\clearpage}
.
\documentclass[11pt,a4paper]{article}
\usepackage[T1]{fontenc}
\usepackage[utf8]{inputenc}
\usepackage{mathpazo}
\usepackage[scaled=0.92]{helvet}
\usepackage{amsmath,amssymb}
\usepackage[margin=2.5cm]{geometry}
\usepackage{microtype}
\usepackage{graphicx}
\usepackage{booktabs,array,tabularx}
\usepackage[font=small,labelfont=bf]{caption}
\usepackage{enumitem}
\usepackage[authoryear,round]{natbib}
\usepackage{titlesec}
\usepackage[titles]{tocloft}
\renewcommand{\cftsecfont}{\bfseries}\renewcommand{\cftsecpagefont}{\bfseries}
\renewcommand{\cftsecleader}{\hfill}\setlength{\cftsecnumwidth}{3.2em}
\usepackage{gpt-status}
\usepackage[hidelinks]{hyperref}
\setlength{\parskip}{4pt}
\setlength{\emergencystretch}{3em}
\setlength{\parindent}{0pt}
\setlist{itemsep=2pt,topsep=3pt}
\titleformat{\section}{\Large\bfseries}{\thesection}{1em}{}
\titleformat{\subsection}{\large\bfseries}{\thesubsection}{1em}{}
\renewcommand{\thesection}{\gptpart.\arabic{section}}
\renewcommand{\thetable}{\gptpart.\arabic{table}}
\renewcommand{\thefigure}{\gptpart.\arabic{figure}}
\renewcommand{\theequation}{\gptpart.\arabic{equation}}
\setcounter{tocdepth}{1}
\makeatletter\renewcommand{\@pnumwidth}{2.2em}\renewcommand{\@tocrmarg}{3em}
\makeatother
\bibliographystyle{plainnat}
\newcommand{\gptitle}{%
  \begin{titlepage}\centering
  \vspace*{3cm}
  {\Huge\bfseries Graviton Pressure Theory\par}\vspace{0.8cm}
  {\Large The Unified Framework\par}\vspace{0.2cm}
  {\large Individual Submission\par}\vspace{0.6cm}
  {\normalsize This document is part of a multi-part scientific framework\par}\vspace{2cm}
  {\normalsize Part \gptpart\ of 30\par}\vspace{0.8cm}
  {\huge\bfseries \gpttitle\par}\vspace{0.6cm}
  {\Large \gptsubtitle\par}\vspace{1.5cm}
  {\small\begin{minipage}{0.75\textwidth}\centering This submission is part of the broader Graviton Pressure Theory (GPT) project, a redefinition of gravitational interaction rooted in causal field dynamics and physical force transmission. Each document is designed to stand on its own; its full context emerges as part of the larger framework. For the complete series, see https://peopleoftheworld.net/graviton-pressure-theory/.\end{minipage}\par}
  \vfill
  {\large Author: Shareef Ali Rashada\par}\vspace{0.2cm}
  {\normalsize shareef@peopleoftheworld.net\par}\vspace{0.5cm}
  {\normalsize Second edition, \gptdate\par}
  \end{titlepage}
  \setcounter{page}{1}\tableofcontents\clearpage}
.
\documentclass[11pt,a4paper]{article}
\usepackage[T1]{fontenc}
\usepackage[utf8]{inputenc}
\usepackage{mathpazo}
\usepackage[scaled=0.92]{helvet}
\usepackage{amsmath,amssymb}
\usepackage[margin=2.5cm]{geometry}
\usepackage{microtype}
\usepackage{graphicx}
\usepackage{booktabs,array,tabularx}
\usepackage[font=small,labelfont=bf]{caption}
\usepackage{enumitem}
\usepackage[authoryear,round]{natbib}
\usepackage{titlesec}
\usepackage[titles]{tocloft}
\renewcommand{\cftsecfont}{\bfseries}\renewcommand{\cftsecpagefont}{\bfseries}
\renewcommand{\cftsecleader}{\hfill}\setlength{\cftsecnumwidth}{3.2em}
\usepackage{gpt-status}
\usepackage[hidelinks]{hyperref}
\setlength{\parskip}{4pt}
\setlength{\emergencystretch}{3em}
\setlength{\parindent}{0pt}
\setlist{itemsep=2pt,topsep=3pt}
\titleformat{\section}{\Large\bfseries}{\thesection}{1em}{}
\titleformat{\subsection}{\large\bfseries}{\thesubsection}{1em}{}
\renewcommand{\thesection}{\gptpart.\arabic{section}}
\renewcommand{\thetable}{\gptpart.\arabic{table}}
\renewcommand{\thefigure}{\gptpart.\arabic{figure}}
\renewcommand{\theequation}{\gptpart.\arabic{equation}}
\setcounter{tocdepth}{1}
\makeatletter\renewcommand{\@pnumwidth}{2.2em}\renewcommand{\@tocrmarg}{3em}
\makeatother
\bibliographystyle{plainnat}
\newcommand{\gptitle}{%
  \begin{titlepage}\centering
  \vspace*{3cm}
  {\Huge\bfseries Graviton Pressure Theory\par}\vspace{0.8cm}
  {\Large The Unified Framework\par}\vspace{0.2cm}
  {\large Individual Submission\par}\vspace{0.6cm}
  {\normalsize This document is part of a multi-part scientific framework\par}\vspace{2cm}
  {\normalsize Part \gptpart\ of 30\par}\vspace{0.8cm}
  {\huge\bfseries \gpttitle\par}\vspace{0.6cm}
  {\Large \gptsubtitle\par}\vspace{1.5cm}
  {\small\begin{minipage}{0.75\textwidth}\centering This submission is part of the broader Graviton Pressure Theory (GPT) project, a redefinition of gravitational interaction rooted in causal field dynamics and physical force transmission. Each document is designed to stand on its own; its full context emerges as part of the larger framework. For the complete series, see https://peopleoftheworld.net/graviton-pressure-theory/.\end{minipage}\par}
  \vfill
  {\large Author: Shareef Ali Rashada\par}\vspace{0.2cm}
  {\normalsize shareef@peopleoftheworld.net\par}\vspace{0.5cm}
  {\normalsize Second edition, \gptdate\par}
  \end{titlepage}
  \setcounter{page}{1}\tableofcontents\clearpage}
.
\documentclass[11pt,a4paper]{article}
\usepackage[T1]{fontenc}
\usepackage[utf8]{inputenc}
\usepackage{mathpazo}
\usepackage[scaled=0.92]{helvet}
\usepackage{amsmath,amssymb}
\usepackage[margin=2.5cm]{geometry}
\usepackage{microtype}
\usepackage{graphicx}
\usepackage{booktabs,array,tabularx}
\usepackage[font=small,labelfont=bf]{caption}
\usepackage{enumitem}
\usepackage[authoryear,round]{natbib}
\usepackage{titlesec}
\usepackage[titles]{tocloft}
\renewcommand{\cftsecfont}{\bfseries}\renewcommand{\cftsecpagefont}{\bfseries}
\renewcommand{\cftsecleader}{\hfill}\setlength{\cftsecnumwidth}{3.2em}
\usepackage{gpt-status}
\usepackage[hidelinks]{hyperref}
\setlength{\parskip}{4pt}
\setlength{\emergencystretch}{3em}
\setlength{\parindent}{0pt}
\setlist{itemsep=2pt,topsep=3pt}
\titleformat{\section}{\Large\bfseries}{\thesection}{1em}{}
\titleformat{\subsection}{\large\bfseries}{\thesubsection}{1em}{}
\renewcommand{\thesection}{\gptpart.\arabic{section}}
\renewcommand{\thetable}{\gptpart.\arabic{table}}
\renewcommand{\thefigure}{\gptpart.\arabic{figure}}
\renewcommand{\theequation}{\gptpart.\arabic{equation}}
\setcounter{tocdepth}{1}
\makeatletter\renewcommand{\@pnumwidth}{2.2em}\renewcommand{\@tocrmarg}{3em}
\makeatother
\bibliographystyle{plainnat}
\newcommand{\gptitle}{%
  \begin{titlepage}\centering
  \vspace*{3cm}
  {\Huge\bfseries Graviton Pressure Theory\par}\vspace{0.8cm}
  {\Large The Unified Framework\par}\vspace{0.2cm}
  {\large Individual Submission\par}\vspace{0.6cm}
  {\normalsize This document is part of a multi-part scientific framework\par}\vspace{2cm}
  {\normalsize Part \gptpart\ of 30\par}\vspace{0.8cm}
  {\huge\bfseries \gpttitle\par}\vspace{0.6cm}
  {\Large \gptsubtitle\par}\vspace{1.5cm}
  {\small\begin{minipage}{0.75\textwidth}\centering This submission is part of the broader Graviton Pressure Theory (GPT) project, a redefinition of gravitational interaction rooted in causal field dynamics and physical force transmission. Each document is designed to stand on its own; its full context emerges as part of the larger framework. For the complete series, see https://peopleoftheworld.net/graviton-pressure-theory/.\end{minipage}\par}
  \vfill
  {\large Author: Shareef Ali Rashada\par}\vspace{0.2cm}
  {\normalsize shareef@peopleoftheworld.net\par}\vspace{0.5cm}
  {\normalsize Second edition, \gptdate\par}
  \end{titlepage}
  \setcounter{page}{1}\tableofcontents\clearpage}
